\documentclass[conference]{IEEEtran}
\IEEEoverridecommandlockouts

\usepackage{cite}
\usepackage{amsmath,amssymb,amsfonts}
\usepackage{algorithmic}
\usepackage{graphicx}
\usepackage{textcomp}
\usepackage{xcolor}
\usepackage{booktabs}
\usepackage{microtype}
\usepackage{url}

\begin{document}

\title{BioDINO-Edge: Biologically-Grounded Multi-Scale Vision Foundation Models for Trustworthy Few-Shot Plant Pathology on Resource-Constrained Hardware}

\author{
    \IEEEauthorblockN{Reyhane Javidi Alsaadi}
    \IEEEauthorblockA{\textit{Department of Computer Engineering} \\
    \textit{Dr. Shariaty Technical University} \\
    Tehran, Iran \\
    reyhane.code@gmail.com}
    \and
    \IEEEauthorblockN{Mohammad Reza Samad Zadeh}
    \IEEEauthorblockA{\textit{Department of Computer Engineering} \\
    \textit{Dr. Shariaty Technical University} \\
    Tehran, Iran}
}

\maketitle

\begin{abstract}
Diagnosing foliar plant diseases early and reliably is central to precision agriculture and robotic crop scouting. Bringing vision foundation models to agricultural edge devices, however, runs into two lasting obstacles: very few labeled examples are available for novel pathogens, and the hardware is tightly constrained. We present \textbf{BioDINO-Edge}, an adaptive dual-foundation framework for accurate few-shot plant pathology on resource-constrained hardware. It combines two complementary foundation models: a lightweight self-supervised patch extractor (\textbf{DINOv2 Small}, 22.1M parameters, 21.1\,ms on CPU) and an organismal vision-language model (\textbf{BioCLIP ViT-B/16}) that captures fine-grained pathological hierarchies. On the expanded PlantVillage 5-shot benchmark covering 11 novel intra-species classes (9,330 test images), BioDINO-Edge reaches a \textbf{98.72\% peak accuracy} and a 30-episode mean of \textbf{97.37\% $\pm$ 0.73\%}. On the in-the-wild PlantDoc benchmark, the model attains an episodic mean of \textbf{83.61\% $\pm$ 11.69\%} under 5-way 5-shot evaluation, exceeding recent multimodal baselines by \textbf{13.00} percentage points.
\end{abstract}

\begin{IEEEkeywords}
Plant Pathology, Few-Shot Learning, Vision Foundation Models, BioCLIP, DINOv2, Fisher-Discriminant Alignment, Adaptive Uncertainty Gating, Edge AI, Precision Agriculture.
\end{IEEEkeywords}

\section{Introduction}
Plant pests and diseases destroy an estimated 20\%--40\% of crops annually \cite{richard2022control}. Traditional diagnosis of plant pathogens heavily relies on visual leaf characteristics such as color, shape, and lesions \cite{khakimov2022traditional}. However, there is a pressing need for more efficient and reliable identification methods, as traditional visual inspection is time-consuming, labor-intensive, and potentially unreliable due to overlapping visual symptoms across different diseases \cite{shoaib2023advanced}.

Deep learning provides an automated alternative for plant disease identification by learning discriminative features directly from images \cite{shoaib2023advanced}. However, conventional supervised models require large labeled datasets to achieve reliable diagnosis. Such datasets are difficult to acquire in real-world agricultural settings due to the limited availability and diversity of field images \cite{shoaib2023advanced, shafik2023systematic}.

To address this data scarcity, Few-Shot Learning (FSL) has emerged as a practical paradigm that enables effective model training from only a few labeled samples. Recent research has attempted to apply FSL using cross-domain few-shot learning models---such as Yang \textit{et al.} (CDFSL-BDC, MAML, and NML) \cite{yang2024laboratory} and Zhao \textit{et al.} (MAFDE-DN4) \cite{zhao2024mafde}---as well as few-shot metric domain adaptation (Kudo \textit{et al.}, FMDA) \cite{kudo2025few} and visual-semantic fusion (Zhou \textit{et al.}, TSMF) \cite{zhou2026vsf}, alongside recent edge-federated approaches (Sharif \textit{et al.}) \cite{sharif2026real}. However, these approaches exhibit several critical limitations:
\begin{itemize}
    \item \textbf{Domain Shift \& Target-Domain Dependency:} Cross-domain models (such as Yang \textit{et al.} \cite{yang2024laboratory} and Zhao \textit{et al.} \cite{zhao2024mafde}) perform poorly under cross-domain evaluation due to complex backgrounds and multi-scale lesions, while metric adaptation methods (Kudo \textit{et al.} \cite{kudo2025few}) still require labeled samples from the target domain.
    \item \textbf{Semantic Misalignment:} Multimodal models relying on LLM text descriptions (e.g., GPT-4o in Zhou \textit{et al.} \cite{zhou2026vsf}) fail to generalize well to real-world examples, suffering from semantic inaccuracies and a lack of fine-grained alignment.
    \item \textbf{Deployment Inefficiencies:} Federated and edge models (such as Sharif \textit{et al.} \cite{sharif2026real}) lack sufficient testing using real-world data streams, making them inefficient for real-time edge deployment.
\end{itemize}

To overcome these challenges, we propose \textbf{BioDINO-Edge}, an adaptive dual-foundation framework designed for real-world, cross-domain few-shot plant pathology on edge devices. BioDINO-Edge combines visual embeddings from DINOv2 \cite{oquab2023dinov2} with clean botanical domain prompts from BioCLIP \cite{stevens2024bioclip}, eliminating the need to fine-tune foundation models to detect subtle necrotic lesions. Using closed-form Fisher-LDA \cite{fisher1936use} and Ledoit--Wolf shrinkage \cite{ledoit2004well}, our model completely removes backpropagation during adaptation, enabling rapid on-device calibration in milliseconds.

The key contributions of this work are as follows:
\begin{enumerate}
    \item \textbf{Dual-Foundation Foliar Representation:} Uniting self-supervised visual embeddings from DINOv2 \cite{oquab2023dinov2} with BioCLIP's \cite{stevens2024bioclip} visual and textual botanical representations to resolve subtle necrotic lesions without foundation-model fine-tuning.
    \item \textbf{Closed-Form Analytical Adaptation:} Integrating Fisher-LDA with Ledoit--Wolf covariance shrinkage and soft-anchored transductive recalibration to overcome high-dimensionality issues and 5-shot centroid shifts within 25 milliseconds without backpropagation.
    \item \textbf{Adaptive Two-Tier Edge Deployment:} A joint margin and Shannon entropy \cite{shannon1948mathematical} gate dynamically routes queries between a standalone mobile app and an on-demand deep module, enabling efficient execution within a constrained storage budget.
    \item \textbf{Rigorous Cross-Domain Validation:} Achieving a mean accuracy of 97.37\% $\pm$ 0.73\% (98.72\% peak) across 9,330 queries on the 11-class 5-shot PlantVillage benchmark and 83.61\% $\pm$ 11.69\% on the in-the-wild PlantDoc benchmark under the standard 5-way 5-shot protocol, outperforming the multimodal VSF baseline \cite{zhou2026vsf} by 13.00 percentage points.
\end{enumerate}

\begin{thebibliography}{14}

\bibitem{richard2022control}
B.~Richard, A.~Qi, and B.~D.~L. Fitt, ``Control of crop diseases through integrated crop management to deliver climate-smart farming systems for low and high input crop production,'' \emph{Plant Pathology}, vol.~71, no.~1, pp. 187--206, 2022.

\bibitem{khakimov2022traditional}
A.~Khakimov, I.~Salakhutdinov, A.~Omolikov, and S.~Utaganov, ``Traditional and current-prospective methods of agricultural plant diseases detection: A review,'' in \emph{IOP Conference Series: Earth and Environmental Science}, vol. 951, no.~1, 2022, p. 012002.

\bibitem{shoaib2023advanced}
M.~Shoaib, B.~Shah, S.~El-Sappagh, A.~Ali, A.~Ullah, F.~Alenezi, T.~Gechev, T.~Hussain, and F.~Ali, ``An advanced deep learning models-based plant disease detection: A review of recent research,'' \emph{Frontiers in Plant Science}, vol.~14, p. 1158933, 2023.

\bibitem{shafik2023systematic}
W.~Shafik, A.~Tufail, A.~Namoun, L.~De~Silva, and R.~Apong, ``A systematic literature review on plant disease detection: motivations, classification techniques, datasets, challenges, and future trends,'' \emph{IEEE Access}, vol.~11, pp. 59\,174--59\,203, 2023.

\bibitem{yang2024laboratory}
S.~Yang, Q.~Feng, J.~Zhang, W.~Yang, W.~Zhou, and W.~Yan, ``From laboratory to field: cross-domain few-shot learning for crop disease identification in the field,'' \emph{Frontiers in Plant Science}, vol.~15, p. 1434222, 2024.

\bibitem{zhao2024mafde}
Y.~Zhao, Z.~Zhang, N.~Wu, Z.~Zhang, and X.~Xu, ``MAFDE-DN4: Improved few-shot plant disease classification method based on deep nearest neighbor neural network,'' \emph{Computers and Electronics in Agriculture}, vol. 226, p. 109373, 2024.

\bibitem{kudo2025few}
S.~Kudo, S.~Kagiwada, and H.~Iyatomi, ``Few-shot metric domain adaptation: Practical learning strategies for an automated plant disease diagnosis,'' \emph{arXiv preprint arXiv:2412.18859}, 2024.

\bibitem{zhou2026vsf}
H.~Zhou, X.~Li, Y.~Nie, H.~Yu, H.~Hu, and L.~Lin, ``VSF: A multimodal few-shot plant disease recognition method in agricultural scenarios,'' \emph{AgriEngineering}, vol.~8, no.~10, p. 407, 2026.

\bibitem{sharif2026real}
M.~I. Sharif, Y.~Zhong, M.~Z. Sajid, and F.~Marinello, ``Real-field-ready and digitally sustainable plant disease recognition via federated multimodal edge learning and few-shot domain adaptation,'' \emph{Agriculture}, vol.~16, no.~9, p. 918, 2026.

\bibitem{oquab2023dinov2}
M.~Oquab \emph{et al.}, ``DINOv2: Learning robust visual features without supervision,'' \emph{Transactions on Machine Learning Research}, 2024.

\bibitem{stevens2024bioclip}
S.~Stevens \emph{et al.}, ``BioCLIP: A vision foundation model for the Tree of Life,'' in \emph{Proc. IEEE/CVF Conf. Computer Vision and Pattern Recognition (CVPR)}, 2024, pp. 19412--19424.

\bibitem{fisher1936use}
R.~A. Fisher, ``The use of multiple measurements in taxonomic problems,'' \emph{Annals of Eugenics}, vol.~7, no.~2, pp. 179--188, 1936.

\bibitem{ledoit2004well}
O.~Ledoit and M.~Wolf, ``A well-conditioned estimator for large-dimensional covariance matrices,'' \emph{Journal of Multivariate Analysis}, vol.~88, no.~2, pp. 365--411, 2004.

\bibitem{shannon1948mathematical}
C.~E. Shannon, ``A mathematical theory of communication,'' \emph{Bell System Technical Journal}, vol.~27, no.~3, pp. 379--423, Jul. 1948; vol.~27, no.~4, pp. 623--656, Oct. 1948.

\end{thebibliography}

\end{document}

